\documentclass[11pt,letterpaper]{article}
\usepackage[T1]{fontenc}
\usepackage[utf8]{inputenc}
\usepackage{lmodern}
\usepackage[margin=1in,headheight=15pt]{geometry}
\usepackage{amsmath,amssymb,amsthm,mathtools}
\usepackage{microtype}
\usepackage[table]{xcolor}
\usepackage{booktabs,array,longtable}
\usepackage{enumitem}
\usepackage{fancyhdr}
\usepackage{titlesec}
\usepackage{hyperref}
\usepackage{bookmark}
\definecolor{ink}{HTML}{17324D}
\definecolor{accent}{HTML}{246B79}
\definecolor{pale}{HTML}{EEF4F6}
\hypersetup{colorlinks=true,linkcolor=ink,citecolor=accent,urlcolor=accent,
 pdftitle={A Number-Valued Broadcast Sum of Surreal Sequences},
 pdfauthor={Research manuscript prepared with ChatGPT}}
\titleformat{\section}{\large\bfseries\color{ink}}{\thesection}{0.65em}{}
\titleformat{\subsection}{\normalsize\bfseries\color{ink}}{\thesubsection}{0.65em}{}
\pagestyle{fancy}
\fancyhf{}
\fancyhead[L]{\small\scshape A broadcast sum of surreal sequences}
\fancyhead[R]{\small Research manuscript}
\fancyfoot[C]{\thepage}
\renewcommand{\headrulewidth}{0.3pt}
\setlength{\parskip}{0.35em}
\setlength{\parindent}{1.2em}
\setlist{nosep,leftmargin=1.8em}
\allowdisplaybreaks[2]
\newtheorem{theorem}{Theorem}[section]
\newtheorem{lemma}[theorem]{Lemma}
\newtheorem{proposition}[theorem]{Proposition}
\newtheorem{corollary}[theorem]{Corollary}
\theoremstyle{definition}
\newtheorem{definition}[theorem]{Definition}
\newtheorem{example}[theorem]{Example}
\theoremstyle{remark}
\newtheorem{remark}[theorem]{Remark}
\newcommand{\No}{\mathbf{No}}
\newcommand{\NN}{\mathbb N}
\newcommand{\RR}{\mathbb R}
\newcommand{\BB}{\mathsf B}
\newcommand{\bs}{\boldsymbol{s}}
\newcommand{\bt}{\boldsymbol{t}}
\newcommand{\bu}{\boldsymbol{u}}
\newcommand{\bv}{\boldsymbol{v}}
\newcommand{\ns}{\mathbin{\#}}
\newcommand{\natsum}{\mathop{\#}}
\newcommand{\len}{\ell}
\newcommand{\cut}[2]{\left\{\,#1\;\middle|\;#2\,\right\}}
\newcommand{\pref}[2]{#1\mathbin{\restriction}#2}
\newcommand{\Rk}{\mathcal R}
\newcommand{\epset}{\mathcal E}
\newcommand{\signplus}{\mathtt{+}}
\newcommand{\signminus}{\mathtt{-}}
\newcommand{\finite}{\mathrm{fin}}
\newcommand{\Ord}{\mathrm{Ord}}

\begin{document}
\hypersetup{pageanchor=false}
\begin{titlepage}
\centering
\vspace*{0.5in}
{\small\scshape\color{accent} Ordinal arithmetic \; / \; Surreal numbers \; / \; Infinite games\par}
\vspace{0.45in}
{\LARGE\bfseries\color{ink} A Number-Valued Broadcast Sum\\[0.18em]
of Surreal Sequences\par}
\vspace{0.25in}
{\large Commuting truncations and an exact failure of monotonicity\par}
\vspace{0.3in}
{\normalsize A partial attack on Lipparini's Problem 7.5\par}
\vspace{0.35in}
{\normalsize Research manuscript prepared with ChatGPT\par}
{\normalsize September 20, 2026\par}
\vspace{0.45in}
\begin{minipage}{0.91\textwidth}
\small
\textbf{Abstract.}
We analyze the sign-truncation rule proposed by Lipparini for extending an
infinitary ordinal operation to surreal sequences. A player chooses a sign
position and applies the associated truncation to all but finitely many
coordinates. We supply an explicit descending ordinal rank and prove that every
resulting game is a number: opposite truncations commute, and the earlier
selected sign survives the other move. The operation agrees with ordinary
surreal addition on finite-support families and with Lipparini's operation on
ordinal sequences. A finite-birthday perturbation bound permits exact
calculations for cofinite dyadic families. The main example is a positive,
classically convergent series with terms either $2^{-(n+d)}$ or
$3\cdot2^{-(n+d+1)}$. Its broadcast value is $k+1$ if exactly $k$ terms take the
second value, and $\sqrt{\omega}$ if infinitely many do. A strictly larger
geometric sequence has broadcast value $1$. Thus the proposed rule is
number-valued but fails even weak coordinatewise monotonicity, including on
absolutely convergent positive dyadic sequences. Full proofs and an exact finite
verification program are provided.
\end{minipage}
\vfill
\begin{minipage}{0.91\textwidth}
\colorbox{pale}{\parbox{\dimexpr\linewidth-2\fboxsep\relax}{\small
\textbf{Status and scope.} The underlying rule is Lipparini's, not a new
definition claimed here. The proofs below are research claims developed in this
manuscript; they have not been independently refereed or formally verified.
The result analyzes a published candidate for an open-ended extension problem.
It does \emph{not} construct a weakly monotone surreal summation theory or prove
that no such alternative is possible. A limited literature check does not
establish priority.}}
\end{minipage}
\vspace*{0.3in}
\end{titlepage}
\pagenumbering{roman}
\hypersetup{pageanchor=true}
\setcounter{tocdepth}{1}
\tableofcontents
\vspace{1.3em}
\noindent\textbf{Reading guide.} Sections~1--3 state the target and the rule.
Sections~4--7 establish the structural theorems. Sections~8--10 contain the
exact dyadic calculation and counterexamples. Sections~11--13 and the appendices
document verification, limitations, and reproducibility.
\newpage
\pagenumbering{arabic}

\section{The target and the outcome}\label{sec:target}

Lipparini's extended paper \cite{lipparini2026} asks in Problem~7.5 for a suitable
surreal or game-theoretic extension of the ordinal game in Definition~7.1.
Remark~7.6(3) proposes the sign-truncation rule studied here and already observes
that its runs terminate. The version checked is arXiv:2505.00424v2, dated
April~30, 2026; the relevant discussion is on printed pages~36--37.
We use the local name \emph{broadcast sum} and the notation $\BB$ for this rule.

There are two distinct questions. Does the rule yield a surreal number rather
than an arbitrary game? And does it retain useful order or summation properties?
Our answer is positive to the first question and sharply negative to the second.

\subsection{Main results in a compact form}

For every set-indexed family $\bs=(s_i)_{i\in I}$ of canonical surreal sign
expansions, the rule defines a well-founded game $G(\bs)$. We prove
\begin{equation}\label{eq:main-number}
                     G(\bs)\text{ is a number},\qquad \BB(\bs)\in\No.
\end{equation}
The passage from countable to set-indexed families uses the same finite-exemption
rule; no cardinal-dependent normalization is introduced.

For a concrete exact calculation, let $d\geq2$ be an integer and $E\subseteq\NN$,
where $\NN=\{0,1,2,\ldots\}$. Put
\begin{equation}\label{eq:main-family}
 s_n^{E,d}=\begin{cases}
 3\cdot2^{-(n+d+1)},&n\in E,\\
 2^{-(n+d)},&n\notin E.
 \end{cases}
\end{equation}
Then
\begin{equation}\label{eq:main-formula}
 \boxed{\quad
 \BB\bigl((s_n^{E,d})_{n<\omega}\bigr)=
 \begin{cases}
 |E|+1,&E\text{ finite},\\[2pt]
 \sqrt{\omega},&E\text{ infinite}.
 \end{cases}\quad}
\end{equation}
Here $\sqrt{\omega}$ is the positive square root in the surreal field. It is not
ordinal exponentiation by a fractional exponent. In Conway monomial notation
it is $\omega^{1/2}$.

Every series in \eqref{eq:main-family} is an ordinary positive real series, and
\begin{equation}\label{eq:classical-bound}
 0<\sum_{n=0}^{\infty}s_n^{E,d}\leq3\cdot2^{-d}<1.
\end{equation}
Nevertheless, the broadcast output can exceed every real number. More strongly,
$s_n^{E,d}<2^{-(n+d-1)}$ for every $n$, whereas
\begin{equation}\label{eq:larger-sequence}
 \BB\bigl((2^{-(n+d-1)})_{n<\omega}\bigr)=1.
\end{equation}
Thus \eqref{eq:main-formula} is an explicit counterexample to weak monotonicity
whenever $E\neq\varnothing$.

\subsection{What is and is not resolved}

The manuscript supplies a numerical extension of the ordinal rule and a
substantial analysis of that extension. It also rules out an attractive property
of this particular candidate. This is a partial resolution of the stated
research direction, not a claim that the word ``suitable'' has acquired a unique
mathematical meaning.

The strongest new claims proposed here are the general number-valuedness theorem,
the finite-short-perturbation estimate, and the exact classification
\eqref{eq:main-formula}. The explicit rank is a quantitative proof of termination,
not a claim to originate the termination observation. Standard background on
numbers, cuts, and game comparison is used throughout
\cite{conway,gonshor,schleicherstoll}. Other infinitary game constructions, such
as \cite{lipparini2024,lipparini2025}, should not be identified with the rule
analyzed here.

\section{Conventions and elementary surreal background}\label{sec:background}

\subsection{Signs, prefixes, and options}

We work in ZFC. Index families are set-sized; $\No$ is the proper class of
surreal numbers, used through set-sized sign expansions and cuts.

A surreal number has a unique canonical sign expansion
$s:\len(s)\longrightarrow\{\signplus,\signminus\}$, with ordinal length
$\len(s)$. Its canonical left options are prefixes immediately preceding a plus
sign; its canonical right options are prefixes immediately preceding a minus
sign:
\begin{equation}\label{eq:signcut}
 s=\cut{\pref{s}{\gamma}:s(\gamma)=\signplus}
        {\pref{s}{\gamma}:s(\gamma)=\signminus}.
\end{equation}
Only positions $\gamma<\len(s)$ occur here. A prefix preceding a plus is strictly
smaller than $s$, and a prefix preceding a minus is strictly larger. The empty
string is $0$. The all-plus string of length $\alpha$ represents the ordinal
$\alpha$.

A \emph{short} number in this paper means a number of finite birthday, or
equivalently a dyadic rational represented by a finite sign string. In
particular, ``short'' does not mean every real number. A \emph{finite} surreal
means one whose absolute value is bounded by some natural number; such a number
need not be real or dyadic.

We use game forms when discussing options and their numerical equivalence
classes when calculating values. A game is recursively a \emph{number} when its
options are numbers and every left option is strictly smaller than every right
option. This recursive requirement matters: separation of options without their
number-valuedness is not enough.

\subsection{Comparison and cofinal replacement of cuts}

For numerical game forms $x$ and $y$, the usual comparison test is
\begin{equation}\label{eq:comparison}
 x\leq y
 \quad\Longleftrightarrow\quad
 (\forall x^L\;x^L<y)\ \text{and}\ (\forall y^R\;x<y^R).
\end{equation}
A separated cut denotes the simplest surreal strictly between its two sides.
Replacing a left side by a mutually cofinal set, or a right side by a mutually
coinitial set, leaves the value unchanged. For example, any unbounded set of
nonnegative real numbers is mutually cofinal with $\NN$.

When $L$ and $R$ are sets of numbers, ``mutually coinitial'' means that each
member of either set lies above or equals a member of the other. This is an
order statement, not the assertion that either set has an infimum in $\No$.
These basic tools follow from the ordinary theory of numbers and the simplicity
theorem; see, for instance, \cite[Sections~3--5]{schleicherstoll}.

\subsection{Arithmetic conventions}

We write $\alpha\ns\beta$ for the finite Hessenberg natural sum of ordinals.
In common Cantor normal form it adds the corresponding finite coefficients;
it is strictly increasing in either argument. The natural sum of a finite
family includes every indexed occurrence, with multiplicity.

In Section~\ref{sec:rank}, addition, multiplication, and powers of ordinals
without $\ns$ mean \emph{ordinary ordinal} operations. In the remaining sections, products such as $\omega 2^{-m}$ and
expressions such as $\omega-k$ mean \emph{surreal field} arithmetic.
This distinction is essential: $\omega/2$ and $\omega-1$ are surreals, not
ordinals.

For later use, the standard Conway omega-map satisfies
\begin{equation}\label{eq:omega-map}
 \omega^x=
 \cut{0,\ n\omega^{x^L}:n\in\NN_{>0}}
      {2^{-n}\omega^{x^R}:n\in\NN},
 \qquad \omega^{x+y}=\omega^x\omega^y.
\end{equation}
Using $1/2=\cut{0}{1}$ gives the particular standard identity
\begin{equation}\label{eq:sqrt-cut}
 \cut{\NN}{\omega/2^m:m\in\NN}=\omega^{1/2}=\sqrt{\omega}.
\end{equation}
Formula \eqref{eq:omega-map} is a background construction, not an infinite
summation convention introduced here \cite{conway,gonshor}.

\section{The broadcast rule}\label{sec:rule}

\begin{definition}[Truncation maps]\label{def:truncation}
For a sign $\sigma\in\{\signplus,\signminus\}$ and an ordinal $\gamma$, set
\begin{equation}\label{eq:single-truncate}
 T^\sigma_\gamma(s)=
 \begin{cases}
 \pref{s}{\delta},&
 \delta=\min\{\eta:\gamma\leq\eta<\len(s),\ s(\eta)=\sigma\}
 \text{ exists},\\
 s,&\text{otherwise}.
 \end{cases}
\end{equation}
For a finite set $F\subseteq I$, define
\begin{equation}\label{eq:family-truncate}
 (T^\sigma_{\gamma,F}\bs)_i=
 \begin{cases}s_i,&i\in F,\\T^\sigma_\gamma(s_i),&i\notin F.\end{cases}
\end{equation}
The action is \emph{legal} at $\bs$ when some $j\notin F$ has
$\gamma<\len(s_j)$ and $s_j(\gamma)=\sigma$.
\end{definition}

A legal plus action is a Left move and a legal minus action is a Right move.
The associated recursive game is
\begin{equation}\label{eq:game-definition}
 G(\bs)=\cut{G(T^+_{\gamma,F}\bs):\text{legal}}
                  {G(T^-_{\eta,H}\bs):\text{legal}}.
\end{equation}
This reformulation of Lipparini's rule \cite[Remark~7.6(3)]{lipparini2026}
separates the globally defined truncation map from the requirement that a move
have a witness at its exact threshold. That distinction will prevent a subtle
mistake in the proof of number-valuedness.

For a fixed input, all reachable coordinates are prefixes of its original
coordinates. The set of reachable states is consequently contained in a
set-sized product of prefix sets, and the options at every state form sets.
Section~\ref{sec:rank} establishes the well-foundedness needed to justify
\eqref{eq:game-definition}; no proper-class cut is used.

\section{A descending ordinal rank}\label{sec:rank}

Fix an initial family and put
\begin{equation}\label{eq:rank-bound-parameters}
 \Lambda=\sup_{i\in I}\len(s_i),\qquad \Omega=\omega^{\Lambda+1}.
\end{equation}
These parameters remain fixed while considering followers of the initial
family. For any current coordinate $t_i$, define its plus and minus heights by
\begin{equation}\label{eq:polarity-heights}
 p_i=\sup\{\eta+1:t_i(\eta)=\signplus\},\qquad
 m_i=\sup\{\eta+1:t_i(\eta)=\signminus\},
\end{equation}
with empty supremum $0$. Both heights are at most $\Lambda$.

For any ordinal family $\bu=(u_i)_{i\in I}$ bounded by $\Lambda$, set
\begin{align}
 \zeta(\bu)&=\min\{\xi:\{i:u_i>\xi\}\text{ is finite}\},\label{eq:zeta}\\
 h(\bu)&=\natsum_{i:u_i>\zeta(\bu)}u_i,\label{eq:exception-height}\\
 R(\bu)&=\Omega\cdot\zeta(\bu)+h(\bu).\label{eq:one-rank}
\end{align}
The minimum exists since $\xi=\Lambda$ qualifies. The sum in
\eqref{eq:exception-height} is finite. Moreover $h(\bu)<\Omega$, since its
summands are below the pure power $\Omega$ and finite natural sums preserve that
bound. The form of $\zeta$ is inspired by the ordinal threshold invariant used
in \cite{lipparini2026}; our rank tracks both polarities.

\begin{lemma}[Monotonicity of the auxiliary rank]\label{lem:rank-monotone}
If $v_i\leq u_i$ for all $i$, then $R(\bv)\leq R(\bu)$.
If additionally $\zeta(\bv)=\zeta(\bu)$ and at least one coordinate above this
common threshold strictly decreases, then $R(\bv)<R(\bu)$.
\end{lemma}
\begin{proof}
Coordinatewise decrease implies $\zeta(\bv)\leq\zeta(\bu)$. If the inequality is
strict, then
\[
 R(\bv)<\Omega\bigl(\zeta(\bv)+1\bigr)
 \leq\Omega\zeta(\bu)\leq R(\bu).
\]
If the thresholds agree, every new exceptional coordinate was already
exceptional. Write both finite natural sums over the old exceptional index set,
using zero for a coordinate that is no longer exceptional. The assertion follows
from monotonicity, or strict monotonicity, of finite natural sum.
\end{proof}

\begin{theorem}[Well-foundedness and a birthday bound]\label{thm:rank}
For the current family $\bt$, let $\boldsymbol p,\boldsymbol m$ be as in
\eqref{eq:polarity-heights}, and define
\begin{equation}\label{eq:total-rank}
 \Rk(\bt)=R(\boldsymbol p)\ns R(\boldsymbol m).
\end{equation}
Every legal broadcast move strictly decreases $\Rk$. Consequently
$G(\bs)$ is well-defined, no run is infinite, and its game-form birthday is at
most $\Rk(\bs)<\Omega^2$.
\end{theorem}
\begin{proof}
Every truncation replaces strings by prefixes, so both polarity-height families
decrease coordinatewise. Both auxiliary ranks therefore weakly decrease by
Lemma~\ref{lem:rank-monotone}.

Consider a legal Left move at threshold $\gamma$, with witness $j\notin F$.
After the move, every unprotected coordinate has plus height at most $\gamma$:
there are no plus signs between $\gamma$ and its cutting position, and all later
signs are deleted. Thus $\zeta(\boldsymbol p')\leq\gamma$. The witness satisfies
$p_j>\gamma$ and $p'_j\leq\gamma<p_j$.

If the plus threshold $\zeta$ decreases, its rank strictly decreases. Otherwise
both thresholds agree and are at most $\gamma$. The witness was an exceptional
coordinate and strictly decreased, so the second assertion of
Lemma~\ref{lem:rank-monotone} again gives $R(\boldsymbol p')<R(\boldsymbol p)$.
The minus rank does not increase. Strict monotonicity of natural sum gives a
strict decrease of \eqref{eq:total-rank}. The Right case exchanges the two signs.

All reachable states form a set, and their ranks are ordinals. Transfinite
recursion on this rank therefore defines the game. Induction also gives
$b(G(\bt))\leq\Rk(\bt)$: each follower has rank strictly smaller than the current
rank, so its birthday plus one is at most the current rank.

Finally, $\Lambda+1<\Omega$ and
$R(\bu)<\Omega(\Lambda+1)<\Omega^2$. The ordinal $\Omega^2$ is again a pure power
of $\omega$, hence closed under finite natural sums of smaller ordinals. Thus
$\Rk(\bt)<\Omega^2$.
\end{proof}

\begin{remark}
There need not be a uniform finite bound on the length of runs. Termination
means that each individual run is finite, not that the game tree has finite
height. The rank bound is deliberately generous and is not asserted to be
optimal.
\end{remark}

\section{Commuting truncations force number-valuedness}\label{sec:number}

\begin{lemma}[Commutation]\label{lem:commute}
For fixed signs, thresholds, and finite exemption sets, the truncation maps
commute:
\begin{equation}\label{eq:commutation}
 T^\sigma_{\gamma,F}T^\tau_{\eta,H}\bs
   =T^\tau_{\eta,H}T^\sigma_{\gamma,F}\bs.
\end{equation}
The maps need not both be legal at the states where they are applied.
\end{lemma}
\begin{proof}
Work in one coordinate. If either action exempts that coordinate, the claim is
immediate. Otherwise let $a$ and $b$ be its two cutting positions computed in
the original string, treating an absent cutting position as the full string
length. Applying the two prefix operations in either order gives the prefix of
length $\min(a,b)$. If the first cut erases the sign targeted by the second,
the second operation becomes the identity; this still gives the same minimum.
\end{proof}

\begin{lemma}[A nonidentity action can be legalized]\label{lem:legalize}
If $T^\sigma_{\gamma,F}\bs\neq\bs$, its outcome is a legal move by the
$\sigma$-player, possibly at a threshold larger than $\gamma$.
\end{lemma}
\begin{proof}
Choose
\[
 \delta=\min\{\eta\geq\gamma:\exists i\notin F\;
                [\eta<\len(s_i)\text{ and }s_i(\eta)=\sigma]\}.
\]
This set is nonempty because the action is not the identity. There is a witness
at $\delta$, and no relevant unprotected sign occurs in $[\gamma,\delta)$.
It follows coordinatewise that
$T^\sigma_{\gamma,F}\bs=T^\sigma_{\delta,F}\bs$.
\end{proof}

\begin{lemma}[The earlier witness survives]\label{lem:survive}
Suppose a Left action at $\gamma$ and a Right action at $\eta$ are both legal
at $\bs$. If $\gamma\leq\eta$, the original Left action is still legal after
the Right action. If $\eta\leq\gamma$, the original Right action is still legal
after the Left action.
\end{lemma}
\begin{proof}
In the first case, take the chosen plus witness at position $\gamma$. A Right
cut, unless the coordinate is protected, occurs at a minus position at least
$\eta$. This position is strictly larger than $\gamma$: equality would require
a sign to be both plus and minus. Thus the witness survives. The other case is
symmetric.
\end{proof}

\begin{theorem}[Every broadcast game is a number]\label{thm:number}
For every set-indexed family of surreal numbers in their canonical sign
expansions, $G(\bs)$ is a numerical game. Its surreal value $\BB(\bs)$ is
therefore defined.
\end{theorem}
\begin{proof}
Induct on the rank of Theorem~\ref{thm:rank}. All options are numbers by
induction. It remains to compare an arbitrary Left option and an arbitrary
Right option. Write their states as
\[
 \bs_L=T^+_{\gamma,F}\bs,\qquad
 \bs_R=T^-_{\eta,H}\bs,
\]
and use Lemma~\ref{lem:commute} to define the common state
\[
 \bs_* = T^-_{\eta,H}\bs_L=T^+_{\gamma,F}\bs_R.
\]
If $\gamma\leq\eta$, Lemma~\ref{lem:survive} makes $G(\bs_*)$ a genuine Left
option of $G(\bs_R)$, so
$\BB(\bs_*)<\BB(\bs_R)$. Starting instead from $\bs_L$, the action leading to
$\bs_*$ is either the identity or, by Lemma~\ref{lem:legalize}, a genuine Right
move at an appropriate threshold. Hence
\[
                   \BB(\bs_L)\leq\BB(\bs_*)<\BB(\bs_R).
\]
If $\eta\leq\gamma$, the symmetric argument gives
\[
                   \BB(\bs_L)<\BB(\bs_*)\leq\BB(\bs_R).
\]
Thus every left option is smaller than every right option. Since all options
are numbers, the recursive criterion makes $G(\bs)$ a number.
\end{proof}

\begin{remark}[Why a tempting shortcut would be invalid]
One cannot assume that both original thresholds remain legal after the other
move. A sign may be erased. The proof needs the stronger distinction between a
truncation \emph{map} and a legal \emph{move}: at least one witness survives,
and the other action can be legalized or is the identity. Nor does the proof
assume that $\BB$ is coordinatewise monotone; that property will be disproved.
\end{remark}

The canonical birthday of $\BB(\bs)$ is no greater than the birthday of the
numerical game form $G(\bs)$, so Theorem~\ref{thm:rank} also bounds the birthday
of the resulting surreal number.

\section{Basic identities and the ordinal restriction}\label{sec:basic}

\begin{proposition}[Relabelling, zeros, and negation]\label{prop:basic}
The operation $\BB$ is invariant under bijections of the index set and under
inserting or removing zero coordinates. Also
\begin{equation}\label{eq:oddness}
                         \BB((-s_i)_{i\in I})=-\BB((s_i)_{i\in I}).
\end{equation}
\end{proposition}
\begin{proof}
A bijection carries finite exemption sets and legal witnesses to the
corresponding ones, yielding the same recursive game. Empty strings supply no
witness and are unchanged by every truncation, so zero coordinates can be
ignored, including any of them lying in an exemption set. Negation reverses
every sign, interchanges Left and Right actions, and commutes with prefix
truncation. Induction on rank gives the usual game negation and hence
\eqref{eq:oddness}.
\end{proof}

\begin{theorem}[Finite-support agreement]\label{thm:finite-support}
If $s_i=0$ outside a finite subset $J\subseteq I$, then
\begin{equation}\label{eq:finite-support}
                            \BB(\bs)=\sum_{i\in J}s_i,
\end{equation}
where the right side is ordinary finite surreal addition.
\end{theorem}
\begin{proof}
Induct on the rank of the family. Every follower still has finite support, so
its value is the finite sum of its coordinates by induction.

Every ordinary single-coordinate move in the disjunctive sum is available as a
broadcast move: exempt all the other nonzero coordinates. Conversely, a Left
broadcast move replaces every changed coordinate by one of its left prefixes.
Select its witnessing coordinate $j$. The ordinary move changing only $j$ to
the same prefix is a Left option of the ordinary sum and has value at least as
large as the broadcast follower, because all other coordinate changes weakly
lower the finite sum. The analogous Right comparison reverses the inequality.
Thus the Left option sets are mutually cofinal and the Right option sets
mutually coinitial. Their numerical cuts have the same value.
\end{proof}

\begin{corollary}[Ordinal extension]\label{cor:ordinal}
On countable families of ordinals, $\BB$ equals the infinitary ordinal operation
$\#^S$ characterized by Lipparini's Definition~7.1 and Theorem~7.3.
\end{corollary}
\begin{proof}
For the all-plus sign string of length $\alpha_i$, a legal plus action at
$\gamma$ replaces an unprotected coordinate by $\min(\alpha_i,\gamma)$.
There are no Right moves. This is exactly the ordinal game, whose value is
$\#^S$ by \cite[Theorem~7.3]{lipparini2026}.
\end{proof}

\begin{remark}[Dependence on the index set]
For any infinite set $I$, $\BB((1)_{i\in I})=\omega$: a Left move retains an
arbitrary finite number of ones, and there are no Right moves. Thus the
set-indexed version is not a cardinal-valued counting operation. Its extension
beyond sequences is a consequence of the same rules, not a proposed solution
to every uncountable summation problem.
\end{remark}

\section{A finite-short-perturbation estimate}\label{sec:perturbation}

A useful comparison principle survives despite the later failure of
coordinatewise monotonicity. It controls changes in a finite number of
finite-birthday coordinates.

\begin{theorem}[Finite-short-perturbation bound]\label{thm:perturbation}
Suppose $\bs$ and $\bt$ have the same index set and coincide outside a finite
set $E$. Suppose that $s_i$ and $t_i$ are short for $i\in E$. Define
\begin{equation}\label{eq:perturbation-budget}
                    N=\sum_{i\in E}\bigl(\len(s_i)+\len(t_i)\bigr)\in\NN.
\end{equation}
Then
\begin{equation}\label{eq:perturbation-bound}
                       |\BB(\bs)-\BB(\bt)|\leq N.
\end{equation}
\end{theorem}
\begin{proof}
Use a common ordinal bound $\Lambda$ for both families and the same rank
parameters as in Section~\ref{sec:rank}. Prove
$\BB(\bs)\leq\BB(\bt)+N$ by induction on
$\Rk(\bs)\ns\Rk(\bt)$, keeping the finite set $E$ fixed.

Apply a truncation map, with the same sign, threshold, and exemption set, to
both families, giving $\bs'$ and $\bt'$. These still agree outside $E$, and the
new budget $N'$ is at most $N$. By Lemma~\ref{lem:legalize}, each action is either
a legal move or the identity. If one side changes and the other does not, some
coordinate in $E$ must have lost at least one sign: outside $E$ the inputs and
actions agree. Since all lengths in $E$ are finite,
\begin{equation}\label{eq:budget-drop}
       \text{one action proper and the other idle}\quad\Longrightarrow\quad
                            N'\leq N-1.
\end{equation}
Whenever at least one side changes, the rank of the pair strictly decreases,
so the induction hypothesis applies to the resulting pair.

For a Left move $\bs\to\bs'$, match the action on $\bt$. If the matched action
is a proper Left move, then
\[
 \BB(\bs')\leq\BB(\bt')+N'
           \leq\BB(\bt')+N<\BB(\bt)+N.
\]
If it is idle, \eqref{eq:budget-drop} instead gives
\[
 \BB(\bs')\leq\BB(\bt)+N'\leq\BB(\bt)+N-1<\BB(\bt)+N.
\]
Thus every Left option of $G(\bs)$ lies below $\BB(\bt)+N$.

For a Right move $\bt\to\bt'$, match it on $\bs$. If the matched action is a
proper Right move, then
\[
 \BB(\bs)<\BB(\bs')\leq\BB(\bt')+N'\leq\BB(\bt')+N.
\]
If it is idle, induction and \eqref{eq:budget-drop} give
\[
 \BB(\bs)\leq\BB(\bt')+N'\leq\BB(\bt')+N-1<\BB(\bt')+N.
\]
Use the game form $G(\bt)+N$, with $N$ its nonnegative integer form. Its Right
options are exactly $G(\bt)^R+N$, since $N$ has no Right options. The two groups
of inequalities just proved are precisely the comparison test
\eqref{eq:comparison} for $\BB(\bs)\leq\BB(\bt)+N$. Interchanging the two
families proves the reverse bound.
\end{proof}

\begin{corollary}\label{cor:cofinite}
If a family of short numbers agrees outside a finite set with a constant short
number $r$, then its broadcast value differs from $\BB(r,r,\ldots)$ by a finite
surreal number.
\end{corollary}

The corollary is not an infinitesimal estimate, an integer-error formula, or an
$\ell^1$-continuity theorem. It only says that the error is bounded in absolute
value by a finite integer. The distinction will be enough to identify the
leading scale of all Right options in the main example.

\section{Constant dyadic sequences and finite defects}\label{sec:constants}

For $m\in\NN$, write
\begin{equation}\label{eq:r-X}
                         r_m=2^{-m},\qquad X_m=\omega r_m.
\end{equation}
All products and differences in this section are in the surreal field.

\subsection{Two standard cut identities}

\begin{lemma}\label{lem:X-cuts}
The following cuts hold:
\begin{align}
 X_0&=\cut{\NN}{},\label{eq:X0}\\
 X_m&=\cut{\NN}{X_{m-1}-n:n\in\NN}\qquad(m\geq1),\label{eq:Xm}\\
 X_h-k&=\cut{\NN}{X_h-(k-1)}\qquad(h\geq0,\ k\geq1).
 \label{eq:X-minus-k}
\end{align}
\end{lemma}
\begin{proof}
The first identity is the definition of $\omega$. Since
$r_m=\cut{0}{r_{m-1}}$ for $m\geq1$, the Conway product with
$\omega=\cut{\NN}{}$ gives
\[
 X_m=\cut{nr_m:n\in\NN}{X_{m-1}-nr_m:n\in\NN}.
\]
The left side is mutually cofinal with $\NN$; replacing $nr_m$ by $n$ on the
right is also a mutually coinitial replacement. This gives \eqref{eq:Xm}.

For \eqref{eq:X-minus-k}, use the finite-sum form of $X_h+(-k)$, where
$-k=\cut{} {-(k-1)}$. Its left options are real numbers cofinal with $\NN$.
One right option is $X_h-(k-1)$. If $h\geq1$, the other right options have the
form $X_{h-1}-n-k=2X_h-n-k$, all strictly above $X_h-(k-1)$, since $X_h$
exceeds every real number. They can be discarded. For $h=0$ there are no such
additional right options.
\end{proof}

\begin{theorem}[Constant dyadic values and finite defects]\label{thm:constants}
For every $m\in\NN$,
\begin{equation}\label{eq:constant-value}
                              \BB(r_m,r_m,\ldots)=X_m.
\end{equation}
For $m\geq1$, the family with infinitely many $r_{m-1}$ and exactly $k$ entries
equal to $r_m$ has value
\begin{equation}\label{eq:defect-value}
                              H_{m,k}=X_{m-1}-k.
\end{equation}
The positions of the finitely many exceptional entries do not matter.
\end{theorem}
\begin{proof}
We use induction on $m$. The constant sequence of ones has only Left moves,
which give arbitrary finite collections of ones. Its cut is
$\cut{\NN}{}=\omega$, proving \eqref{eq:constant-value} for $m=0$.

Assume the constant formula for all smaller indices. For this fixed $m\geq1$,
first prove \eqref{eq:defect-value} by induction on $k$. Its base case is the
already established constant value $X_{m-1}$.

The sign expansion of $r_j$ is a plus followed by $j$ minuses. Hence the only
Left threshold of the family defining $H_{m,k}$ is $0$. Such a move leaves
finitely many positive dyadic entries and has their ordinary finite sum as its
value, by Theorem~\ref{thm:finite-support}. Since infinitely many $r_{m-1}$ are
available, these left values are unbounded finite reals and are mutually
cofinal with $\NN$.

A Right move at threshold $m$ can act only on the $k$ exceptional entries.
It changes each unprotected $r_m$ into $r_{m-1}$, producing precisely values
$H_{m,j}$ with $0\leq j<k$. Every such $j$ can be obtained by protecting the
appropriate exceptional entries. By the induction on $k$, the smallest such
Right value is $X_{m-1}-(k-1)$.

A Right move at a threshold $q<m$, necessarily $q\geq1$, produces a family
that is constantly $r_{q-1}$ outside finitely many short exceptions. The
perturbation bound and the outer induction show that its value is
$X_{q-1}+\varepsilon$ for a finite surreal $\varepsilon$. Here
$q-1\leq m-2$, so $X_{q-1}\geq2X_{m-1}$. Such a value is strictly greater than
$X_{m-1}$ plus any prescribed finite integer, and is therefore greater than
$X_{m-1}-(k-1)$. These other Right options impose no stronger constraint.
The cut for $H_{m,k}$ is consequently
\[
                     \cut{\NN}{X_{m-1}-(k-1)}=X_{m-1}-k
\]
by Lemma~\ref{lem:X-cuts}. This finishes the induction on $k$.

Now consider the constant $r_m$ family. Its Left values are again finite reals
cofinal with $\NN$. A Right move at threshold $m$, with exactly $k$ entries
protected, gives $H_{m,k}=X_{m-1}-k$, for every $k\in\NN$. Right moves at
$q<m$ are larger by the same finite-perturbation argument. Hence its cut is
\[
                      \cut{\NN}{X_{m-1}-k:k\in\NN}=X_m.
\]
This completes the outer induction.
\end{proof}

\begin{corollary}[Cofinite dyadic leading scale]\label{cor:dyadic-tail}
If $\bt$ is constantly $r_m$ outside finitely many short exceptions, then there
is an $N\in\NN$ such that
\begin{equation}\label{eq:tail-bound}
                              X_m-N\leq\BB(\bt)\leq X_m+N.
\end{equation}
In particular $\BB(\bt)$ exceeds every real number.
\end{corollary}
\begin{proof}
Combine Theorems~\ref{thm:perturbation} and \ref{thm:constants}. A positive real
multiple of $\omega$, minus a finite number, is still larger than every real.
\end{proof}

\section{An exact classification of convergent dyadic series}\label{sec:classification}

Fix $d\geq2$ and abbreviate $\bs^{E,d}$ from \eqref{eq:main-family} to $\bs^E$
when the shift is understood. Let
\[
                        a_n=2^{-(n+d)},\qquad b_n=\tfrac32a_n.
\]
Their sign expansions are
\begin{equation}\label{eq:ab-signs}
 a_n=\signplus\,\signminus^{\,n+d},\qquad
 b_n=\signplus\,\signminus^{\,n+d}\,\signplus.
\end{equation}
The second plus in $b_n$ is at position $n+d+1$, with positions counted from
zero. Its deletion changes $b_n$ exactly to $a_n$.

\subsection{A complete description of the relevant options}

\begin{lemma}[Left options]\label{lem:left-family}
For any $E\subseteq\NN$, the Left moves from $\bs^E$ have the following form.
A threshold-zero move has finite support and value strictly less than $1$.
A positive-threshold move produces $\bs^{E'}$ for a finite subset $E'\subseteq E$.
If $E$ is finite and nonempty, $|E'|<|E|$, and a move with
$|E'|=|E|-1$ is available. If $E$ is infinite, every finite value of $|E'|$ is
attainable.
\end{lemma}
\begin{proof}
At threshold zero every unprotected positive string is replaced by the empty
string. Only finitely many entries remain, and their sum is less than the total
real sum, which is at most $3\cdot2^{-d}<1$.

A positive Left threshold must be $\gamma=n+d+1$ for a second plus at some
$n\in E$. Coordinates of the base type $a_j$ have no plus sign at or beyond
$\gamma$ and are unchanged. In an upgraded coordinate $b_j$, the extra plus is
removed exactly when it occurs at or after $\gamma$ and the coordinate is
unprotected. Only finitely many second pluses occur before $\gamma$, and only
finitely many indices are protected. Thus finitely many upgrades survive.
The witnessing upgrade is removed, so for finite $E$ the number strictly
falls. Exempting all other upgrades permits exactly one removal.

If $E$ is infinite, choose the threshold of its earliest upgrade. Exempt any
$k$ other upgraded indices. All unprotected upgrades are removed, and precisely
$k$ remain. This realizes every $k\in\NN$.
\end{proof}

\begin{lemma}[Right options]\label{lem:right-family}
Every Right option from $\bs^E$ is constantly $r_{\gamma-1}$ outside finitely
many short exceptions, where $\gamma\geq1$ is its threshold. Consequently its
value is $X_{\gamma-1}+\varepsilon$, with $\varepsilon$ finite, and is larger
than every real number. Every integer threshold $\gamma\geq1$ is legal.
\end{lemma}
\begin{proof}
All minus signs occur at positive finite positions. For a fixed $\gamma\geq1$,
all sufficiently large indices satisfy $n+d\geq\gamma$. Both $a_n$ and $b_n$
have a minus at position $\gamma$, and truncation immediately before that
position gives
\[
                         \signplus\,\signminus^{\,\gamma-1}=r_{\gamma-1}.
\]
Only finitely many earlier coordinates and the finite exemption set can differ
from this tail. All exceptional strings are short. Apply
Corollary~\ref{cor:dyadic-tail}. Since arbitrarily long initial minus blocks are
present, a witness exists at every positive finite threshold.
\end{proof}

\subsection{The finite-upgrade case}

\begin{theorem}\label{thm:finite-E}
If $E$ is finite, then $\BB(\bs^E)=|E|+1$.
\end{theorem}
\begin{proof}
Induct on $k=|E|$. For $k=0$, all Left options are finite-support values less
than $1$, including $0$, and all Right options exceed every real number.
The simplest number between the two sides is $1$: it fits the cut, whereas its
only canonical predecessor $0$ does not. Thus the value is $1$.

For $k>0$, Lemma~\ref{lem:left-family} and induction show that a positive
Left threshold gives a value $j+1$ with $j<k$. Such values are at most $k$, and
$k$ itself is attained. Threshold-zero values remain below $1$. Every Right
option is infinite by Lemma~\ref{lem:right-family}. The simplest number larger
than these Left values and smaller than these Right values is $k+1$: it fits,
and every canonical predecessor of the integer $k+1$ is at most $k$ and fails.
\end{proof}

\subsection{The infinitely-many-upgrades case}

\begin{theorem}\label{thm:infinite-E}
If $E$ is infinite, then $\BB(\bs^E)=\sqrt{\omega}$.
\end{theorem}
\begin{proof}
By Lemma~\ref{lem:left-family} and Theorem~\ref{thm:finite-E}, all Left options
are finite numbers and include every positive integer. Thus the Left option set
is mutually cofinal with $\NN$.

Let $\mathcal Q$ be the set of Right option values. We show that it is mutually
coinitial with $\{X_m:m\in\NN\}$.

First take $q\in\mathcal Q$. By Lemma~\ref{lem:right-family}, for some
$\gamma\geq1$ and some finite integer $N$,
\[
                             q\geq X_{\gamma-1}-N>X_\gamma.
\]
The last inequality holds because $X_{\gamma-1}-X_\gamma=X_\gamma$ is infinite.
So some $X_m$ lies below every given Right value.

Conversely, fix $m\in\NN$. Choose a Right move at threshold $\gamma=m+2$ with
empty exemption set; it is legal by Lemma~\ref{lem:right-family}. Its value
$q_m$ satisfies $q_m\leq X_{m+1}+N_m$ for some finite $N_m$, and therefore
\[
                             q_m<X_m,
\]
since $X_m-X_{m+1}=X_{m+1}$ exceeds every finite integer. Thus some Right value
lies below each prescribed $X_m$.

Replacing both option sets by these mutually cofinal and coinitial sets gives
\[
                 \BB(\bs^E)=\cut{\NN}{X_m:m\in\NN}.
\]
By the standard omega-map cut \eqref{eq:sqrt-cut}, this is
$\omega^{1/2}=\sqrt{\omega}$.
\end{proof}

\begin{remark}[Sparsity does not affect the result]
The proof uses only that $E$ is infinite, not any positive density or bounded-gap
assumption. For example, the squares, powers of two, or any arbitrarily sparse
infinite set give the same output. For finite $E$, the output depends only on
its cardinality, not on the locations or total real mass of the upgrades.
\end{remark}

\section{Consequences: monotonicity, scaling, and limits}\label{sec:consequences}

\begin{lemma}[Pure geometric tails]\label{lem:pure-power}
For every integer $e\geq1$,
\begin{equation}\label{eq:pure-power}
                 \BB\bigl((2^{-(n+e)})_{n<\omega}\bigr)=1.
\end{equation}
\end{lemma}
\begin{proof}
The only Left threshold is zero. Its finite-support values are strictly below
the total real sum $2^{1-e}\leq1$, and include $0$. A Right move at any positive
threshold produces a constant positive dyadic tail with finitely many short
exceptions, hence has infinite value by Corollary~\ref{cor:dyadic-tail}.
The simplest number fitting this cut is $1$.
\end{proof}

\begin{corollary}[Failure of weak monotonicity]\label{cor:nonmonotone}
There are strictly positive absolutely convergent dyadic sequences $\bs$ and
$\bt$ such that $s_n<t_n$ for every $n$ but $\BB(\bs)>\BB(\bt)$.
The reversal can be an arbitrarily large integer or an infinite surreal.
\end{corollary}
\begin{proof}
Set $\bs=\bs^{E,d}$ and $t_n=2^{-(n+d-1)}$. Then
$s_n\leq\tfrac32\,2^{-(n+d)}<2\,2^{-(n+d)}=t_n$. By
Lemma~\ref{lem:pure-power}, $\BB(\bt)=1$. For nonempty finite $E$, the other value
is $|E|+1$; for infinite $E$, it is $\sqrt{\omega}$.
\end{proof}

In the particularly simple all-upgraded case with $d=2$, the reversal reads
\begin{align}
 \BB(3/8,3/16,3/32,\ldots)&=\sqrt{\omega},\label{eq:concrete-smaller}\\
 \BB(1/2,1/4,1/8,\ldots)&=1.\label{eq:concrete-larger}
\end{align}
Every term on the first line is strictly smaller than the corresponding term on
the second. Their classical sums are $3/4$ and $1$ respectively.

\begin{corollary}[Failure of finite-scaling compatibility]\label{cor:scaling}
The operation is not homogeneous under multiplication by $3/2$, even on
positive convergent dyadic sequences. It also need not satisfy
$\BB(\bs+\bs)=2\BB(\bs)$.
\end{corollary}
\begin{proof}
For $a_n=2^{-(n+2)}$, the classification gives $\BB(\boldsymbol a)=1$, while
$\BB((3/2)\boldsymbol a)=\sqrt{\omega}\neq3/2$. Also
$\BB(2\boldsymbol a)=1$ by \eqref{eq:pure-power}, not $2$.
\end{proof}

\subsection{A numerical table with exact surreal outputs}

For $d=2$, write $E_k=\{0,1,\ldots,k-1\}$, including $E_0=\varnothing$.
The ordinary sum is
\begin{equation}\label{eq:finite-E-real-sum}
 \sum_{n=0}^{\infty}s_n^{E,2}
       =\frac12+\sum_{n\in E}2^{-(n+3)},\qquad
 \sum_{n=0}^{\infty}s_n^{E_k,2}=\frac34-2^{-(k+2)}.
\end{equation}
The following values are therefore exact, but the two columns describe
different operations.

\begin{center}
\renewcommand{\arraystretch}{1.16}
\begin{tabular}{lll}
\toprule
Upgrade set or sequence & Classical real sum & Broadcast value\\
\midrule
$E=\varnothing$ & $1/2$ & $1$\\
$E=\{0\}$ & $5/8$ & $2$\\
$E=\{0,1\}$ & $11/16$ & $3$\\
$E=\{0,1,2\}$ & $23/32$ & $4$\\
$E=\NN$ & $3/4$ & $\sqrt{\omega}$\\
$(1/2,1/4,1/8,\ldots)$ & $1$ & $1$\\
\bottomrule
\end{tabular}
\end{center}

\subsection{Why finite truncations cannot establish the infinite answer}

Let $\bs^{[N]}$ denote the family retaining only its first $N$ coordinates and
replacing the rest by zero. The finite-support theorem gives
\[
                      \BB(\bs^{[N]})=\sum_{n=0}^{N-1}s_n.
\]
For the all-upgraded sequence with $d=2$, this is
\begin{equation}\label{eq:partial-sum}
                   \BB(\bs^{[N]})=\frac34(1-2^{-N})<\frac34,
\end{equation}
whereas $\BB(\bs)=\sqrt{\omega}$. Finite game evaluation and subsequent real
limiting would therefore compute the wrong operation.

There is no contradiction. An infinite broadcast move truncates infinitely many
coordinates at once while allowing only finitely many exemptions. Its game tree
is not a numerical limit of the values of finite-support truncations. The
proof of the infinite value relies on its full option sets, especially their
cofinal and coinitial structure.

\subsection{Small real mass does not imply a small output}

For every real $\varepsilon>0$, choose $d$ so large that $3\cdot2^{-d}<\varepsilon$.
The all-upgraded family then has $\ell^1$ norm below $\varepsilon$, but its
broadcast value remains $\sqrt{\omega}$. Thus no bound by a real-valued function
of the ordinary $\ell^1$ norm can hold uniformly for these inputs.

Even a perturbation at a single sufficiently late index has arbitrarily small
real mass but changes the value of the base family from $1$ to $2$.
The finite-short-perturbation theorem is consistent with this behavior: it uses
sign lengths, not numerical size, as its budget.

\section{Computational checks and reproducibility}\label{sec:verification}

The archive includes a standard-library Python implementation using exact
rational arithmetic. It recursively constructs the \emph{actual finite game}
for finite families of finite canonical sign strings, checks separation of
Left and Right option values, and chooses the simplest dyadic in the resulting
cut. It does not insert the finite-support theorem as the evaluation rule.

Running \texttt{python code/verify.py} produced the following test counts:
\begin{center}
\renewcommand{\arraystretch}{1.15}
\begin{tabular}{p{0.75\textwidth}r}
\toprule
Check & Count\\
\midrule
Single-string commutation, sign length at most $8$ & $135{,}156$\\
Finite two-coordinate game values, lengths at most $4$ & $961$\\
Finite three-coordinate game values, lengths at most $3$ & $3{,}375$\\
Legal Left--Right diamonds, including exemption sets & $3{,}378$\\
Input pairs checked for negation, reversal, and zero insertion & $225$\\
Dyadic sign formulas and positive-constant-tail truncations & $14{,}560$\\
\bottomrule
\end{tabular}
\end{center}
All checks passed. The recursive evaluator memoized $4{,}336$ finite states.
The exact report, including the scope limitation, is saved as
\texttt{verification/test\_report.json}. Runtime depends on the environment.

The diamond tests explicitly check both details used in the proof: the earlier
witness survives, and the other action is either idle or a legal move after a
possible threshold adjustment. Merely testing commutation would miss the
second issue.

The file \texttt{verification/partial\_sums.csv} contains exact finite partial
sums for the base and all-upgraded $d=2$ families. Its columns naming the infinite
values are labeled as values \emph{from the theorem}; they are not numerical
outputs of an infinite-surreal evaluator.

\paragraph{What the checks establish.}
They test the implementation and provide finite consistency checks for the local
lemmas and finite-support theorem. They do not prove well-foundedness for
arbitrary ordinal-length strings, the universal number theorem, the
finite-perturbation estimate for long tails, or the value $\sqrt{\omega}$.
Those statements are supported by the written proofs, not by extrapolation from
finite experiments. No proof-assistant formalization is included.

\section{Research assessment and remaining questions}\label{sec:assessment}

\subsection{Why this attack was tractable}

The rule is governed by canonical sign positions rather than arbitrary
comparisons of surreal magnitudes. Prefix truncations have a rigid algebra:
they commute, and an earlier opposite-sign witness survives. That gives a local
mechanism forcing the global game to be a number.

The same sign-based mechanism creates the counterexample. A tiny dyadic
increase from $a_n$ to $b_n$ appends one plus sign at a late position. A single
broadcast move can remove all but finitely many such signs. Thus the
\emph{finite versus infinite} combinatorics of the appended signs matter far
more than the real sizes of the perturbations. The surviving finite cases
supply all positive integers as Left options; constant dyadic tails control the
Right options. The final answer is forced by the cut between these two scales.

This explains why the operation can be a valid numerical extension of an
ordinal game while being unsuitable as an order-preserving completion of
ordinary real summation. Number-valuedness is a game-theoretic property;
coordinatewise monotonicity is an additional assertion about its dependence on
input values. One does not imply the other.

\subsection{Precisely what remains open in this manuscript}

An alternative rule might preserve weak coordinatewise monotonicity while
extending the ordinal operation. Nothing proved here excludes that possibility.
One would have to state the desired compatibility axioms explicitly; the
particular broadcast rule already fails both classical convergence agreement
and even elementary finite-scaling identities.

Several narrower directions also remain untreated. The exact value of
$\BB(r,r,\ldots)$ for a general positive surreal $r$ is not established here;
Theorem~\ref{thm:constants} covers the dyadics $2^{-m}$, not all numbers.
Our finite-short-perturbation estimate does not cover arbitrary finite changes
in long surreal coordinates. The ordinal birthday bound is not sharp. And the
classification of positive dyadic sequences with more than one extra sign per
coordinate may exhibit additional scales, but no formula for them is claimed.

A useful continuation would be to characterize a substantial class of inputs
on which the broadcast operation is monotone, or to determine how a controlled
hierarchy of late sign patterns changes its output cut. These are concrete
questions suggested by the proofs rather than claims that the broad extension
problem has been settled.

\subsection{Provenance and priority limits}

The original problem and the exact proposed rule were checked directly against
the extended source, including the rendered page containing the sign rule.
The arXiv record identifies v2 as an extended version of the published article.
Related game-series papers were also identified, but this was not an exhaustive
bibliographic review. In particular, the present document does not assert that
no independent proof or overlapping result exists elsewhere.

The construction belongs to Lipparini. The standard facts about surreal cuts,
comparison, finite arithmetic, and the omega-map belong to the established
literature. The explicit arguments and classification in this manuscript are
submitted as mathematical research claims to be checked, not presented as
results already certified by those sources.

\section{Conclusion}

The proposed sign-truncation rule is more structured than a general infinite
game: it always produces a surreal number. A well-founded ordinal rank makes
the recursion precise, and a commuting-move argument separates every pair of
opposite options. Finite-support agreement and an explicit perturbation bound
then permit exact transfinite calculations.

The decisive calculation is
\[
 \BB\left(\left(2^{-(n+d)}
       \left(1+\tfrac12\mathbf 1_E(n)\right)\right)_{n<\omega}\right)
  =\begin{cases}|E|+1,&E\text{ finite},\\\sqrt{\omega},&E\text{ infinite}.
    \end{cases}
\]
All these inputs have small convergent positive real sums. A strictly larger
geometric input can have broadcast value $1$. This is both positive progress
on the number-valued extension question and an exact obstruction to using this
particular candidate as a weakly monotone summation theory.

\appendix
\section{Proof dependency audit}\label{app:dependencies}

The main arguments have the following noncircular dependency structure.
\begin{center}
\renewcommand{\arraystretch}{1.14}
\begin{tabular}{p{0.36\textwidth}p{0.54\textwidth}}
\toprule
Result & Inputs actually used\\
\midrule
Descending rank & Prefix lengths; ordinal well-ordering; finite natural sum\\
Number-valuedness & Rank induction; commutation; legalizing nonidle maps;
 surviving earlier witnesses; numerical option comparison\\
Finite-support agreement & Number-valuedness; rank induction; monotonicity of
 \emph{finite ordinary} surreal sums\\
Finite-short perturbation & Number-valuedness; pair-rank induction; finite
 birthday budget; game comparison\\
Constant dyadic values & Finite-support agreement; perturbation; ordinary
 product cuts; nested finite induction\\
Finite-upgrade classification & Option descriptions; constant-tail estimates;
 simplicity of natural numbers\\
Infinite-upgrade classification & Finite-upgrade values; constant-tail estimates;
 mutual coinitiality; the standard omega-map cut\\
Failure of monotonicity & Exact classification; separate pure-power calculation\\
\bottomrule
\end{tabular}
\end{center}
In particular, no step uses coordinatewise monotonicity of $\BB$. The only
monotonicity of numerical addition used in the finite-support proof is the
standard monotonicity of an already defined \emph{finite} sum.

\section{A small worked sign example}

For $d=2$, the first few base and upgraded strings are
\[
\begin{array}{c|c|c|c|c}
 n&a_n&\text{signs of }a_n&b_n&\text{signs of }b_n\\\hline
 0&1/4&\mathtt{+--}&3/8&\mathtt{+--+}\\
 1&1/8&\mathtt{+---}&3/16&\mathtt{+---+}\\
 2&1/16&\mathtt{+----}&3/32&\mathtt{+----+}
\end{array}
\]
Consider the infinite all-upgraded family. A Left move at position $3$ with no
exemptions deletes every extra plus: all their positions are at least $3$.
The result is the base family, whose value is $1$. Protecting any $k$ upgraded
coordinates other than the selected witness leaves exactly $k$ upgrades and
produces value $k+1$.

A Right move at position $2$ with no exemptions changes every entry to the
prefix $\mathtt{+-}=1/2$, giving the constant family of value $\omega/2$.
A Right move at a much later position leaves finitely many early exceptions and
produces a much smaller positive dyadic constant tail. Its value is a smaller
positive dyadic multiple of $\omega$, up to a finite error. These observations
are the concrete form of the two option-set arguments in
Theorem~\ref{thm:infinite-E}.

\section{Archive contents and build instructions}

The archive is self-contained apart from standard Python and a LaTeX
installation. It includes the source \texttt{article.tex}, compiled
\texttt{article.pdf}, bibliography \texttt{references.bib}, a \texttt{build.sh}
script, a \texttt{README.md}, two Python modules under \texttt{code/}, and exact
verification artifacts under \texttt{verification/}. No third-party paper or
font file is redistributed.

Run the finite checks with
\begin{verbatim}
python code/verify.py
\end{verbatim}
Build the paper with
\begin{verbatim}
sh build.sh
\end{verbatim}
The latter runs \texttt{pdflatex}, \texttt{bibtex} (or \texttt{bibtex8}), and two
further \texttt{pdflatex} passes. The bibliography records source identifiers and links;
links were checked during preparation on September~20, 2026. A rebuild of the
PDF is not a verification of the mathematical proofs.

\clearpage
\phantomsection
\addcontentsline{toc}{section}{References}
\begingroup
\raggedright
\bibliographystyle{plain}
\bibliography{references}
\endgroup
\end{document}
